\section{Analysis}

\subsection{Designs in GA}
\begin{wraptable}[7]{r}{0.45\textwidth}  
\vspace{-0.2in}
\small
\begin{tabular}{l||ccc}
\toprule
\textbf{Strategy}  & \textbf{SST-5} & \textbf{ASSET} & \textbf{Avg.}\\ \midrule
random  & 48.67\scriptnumber{0.97} & 46.32\scriptnumber{0.32} & 47.50 \\
tournament & 49.70\scriptnumber{0.60} & 46.29\scriptnumber{0.18} & 48.00 \\
{wheel} & \textbf{49.91}\scriptnumber{0.61} & \textbf{46.43}\scriptnumber{0.19} & \textbf{48.17} \\

\bottomrule
\end{tabular}
\caption{Designs in \modelname (GA).}
\label{exp:ana-design-ga}
\end{wraptable}



For \modelname (GA), we apply the roulette wheel selection strategy by default to select parental prompts, contributing to the offspring. 
To further explore the effect of various selection strategies, we compare our approach with another two popular strategies, i.e., tournament~\citep{enwiki:1160627612} and random selection, as presented in Table~\ref{exp:ana-design-ga}. We observe that \modelname (GA) with roulette wheel achieves higher scores, showcasing the effectiveness of this selection method.

\subsection{Designs in DE}
For \modelname (DE), we delve into two key design considerations in adapting the evolutionary operators of DE to discrete prompts: \textbf{1)} mutation on different parts, and \textbf{2)} choosing the current top-performing prompt as ``Prompt 3'' in Figure~\ref{fig:de_cot}. We assess the impact of these design choices on two datasets: Subj, an understanding dataset where \modelname (DE) outperforms \modelname (GA), and ASSET, a generation dataset where both variants demonstrate similar performance.


\begin{wraptable}[8]{r}{0.49\textwidth}  
\small
\centering
\vspace{-0.15in}
\begin{tabular}{ll||ll}
\toprule
\textbf{Mutation} & \textbf{Prompt} 3 & \textbf{Subj} & \textbf{ASSET} \\ \midrule
{Diff} & {best} & \textbf{75.55}\scriptnumber{2.26} & \textbf{46.21}\scriptnumber{0.27} \\
All & best & 69.87\scriptnumber{0.82} &45.73\scriptnumber{0.45}  \\ 
Diff & random & 69.82\scriptnumber{2.47}& 45.89\scriptnumber{0.37} \\
Diff & eliminate & 69.07\scriptnumber{4.21} & 45.90\scriptnumber{0.23} \\
\bottomrule
\end{tabular}
\caption{Designs in \modelname (DE).}
\label{exp:ana-design-de}
\end{wraptable}
\paragraph{Mutation on Different Parts}
To illustrate the benefits of mutating only the different parts, we replace the first two steps in Figure~\ref{fig:de_cot} with the instruction ``Randomly mutate Prompt 1 and Prompt 2'' to allow mutation on all contents in Prompts 1 and 2, denoted as ``All'' in Table~\ref{exp:ana-design-de}. Meanwhile, the original design in \modelname, which mutates only the different parts, is denoted as ``Diff''. 
As shown in Table~\ref{exp:ana-design-de}, the design of mutation on only the different parts consistently yields performance gains across two tasks.
\paragraph{Selection of Prompt 3}
Applying one of the variants of the DE algorithm, in \modelname (DE), we pick the best prompt in the current population as Prompt 3 in Figure~\ref{fig:de_cot}. We validate this design via the following settings: \textbf{1)} Prompt 3 is randomly sampled from the current population, denoted as ``random'' in Table~\ref{exp:ana-design-de}; \textbf{2)} Eliminate the use of Prompt 3 by letting the Basic Prompt directly cross over with the mutated different parts (i.e., remove Step 3 in Figure~\ref{fig:de_cot}), denoted as ``eliminate'' in Tabel~\ref{exp:ana-design-de}. Table~\ref{exp:ana-design-de} clearly demonstrates the importance of introducing Prompt 3. Moreover, it is shown that choosing the best prompt as Prompt 3 is more effective than random sampling.





\subsection{Population Initialization}



\begin{wraptable}[13]{r}{0.45\textwidth}  
\vspace{-0.15in}  
\small  
\centering  
\begin{tabular}{l||cc}  
\toprule  
\textbf{Initialization} & \textbf{GA} & \textbf{DE} \\ 
\midrule
bottom-10 & 47.80\scriptnumber{0.92} & 48.64\scriptnumber{0.15} \\ \midrule  
random-10 & 49.34\scriptnumber{0.53} & \textbf{50.03}\scriptnumber{1.08} \\  
random-5 + var-5 & 49.84\scriptnumber{1.49} & 49.53\scriptnumber{1.04} \\ \midrule  
top-10 & 49.62\scriptnumber{1.00} & 49.61\scriptnumber{2.30} \\  
{top-5 + var-5} & \textbf{49.91}\scriptnumber{0.61} &  49.89\scriptnumber{1.73} \\  
\bottomrule  
\end{tabular}  
\caption{Ablations of the initial population on SST-5, where top-$n$, random-$n$, bottom-$n$ denotes the top-performing, randomly selected, bottom-performing n prompts, and var-$n$ denotes the number of generated $n$ variations.}  
\label{tab:abs-init}  
\end{wraptable}  

We investigate the effect of initial population quality on \modelname. 
We conduct pilot experiments to sort the prompts (designed manually or generated by GPT-3.5) according to their performance on the dev set. We then select bottom, random and top prompts along with their corresponding variations as initial prompts. These variations are generated using the resampling template designed in~\cite{zhou2022large}, shown in Figure~\ref{fig:app-ape_temp} in the Appendix~\ref{app:temp}, which is used to introduce randomness to the initialization. 

Table~\ref{tab:abs-init} demonstrates that: \textbf{1)} Crafted design of initial prompts is not essential, as randomly selecting prompts can achieve a similar performance to selecting the top-performing ones; 
\textbf{2)} When selecting the top-performing prompts, introducing randomness by allowing GPT-3.5 to generate variations can lead to a slight improvement in overall performance; however, when randomly selecting prompts, there is no need to introduce additional randomness for \modelname (DE); 
\textbf{3)} When using top-performing initial prompts, \modelname (GA) performs slightly better than \modelname (DE); however, when starting with bottom-performing initial prompts, \modelname (DE) outperforms \modelname (GA), which indicates that DE is a better choice when the available manual prompts are not of high quality.

 